% Source: verified native ZIP, components/sga4half_intro_en.tex;
% collated against French original PDF pp.3--4 (printed III--IV).
\chapter*{引言}
\addcontentsline{toc}{chapter}{引言}
本卷旨在使非专业读者更容易使用 $\ell$ 进上同调. 我希望它往往能使读者免于查阅 SGA~4 和 SGA~5 中繁密的各讲. 本卷也包含一些新结果.

第一讲由 J.~F. Boutot 撰写, 概览 SGA~4. 它给出主要结果---只采用最低限度的一般性, 对应用而言往往还不够---并说明证明的大意. 要查阅完整的结果或详细的证明, SGA~4 仍不可缺少.

《\nameref{II}》包含 Frobenius 自同态的迹公式的完整证明. 这个证明就是 Grothendieck 在 SGA~5 中给出的证明, 删去了所有不必要的细节. 这篇报告应能使使用者不必再查阅 SGA~5; 后者可以看作一系列离题的讨论, 其中有些十分有趣. 有了这篇报告, SGA~5 就可以很快按原样出版. 《\nameref{VI}》是它的补充, 说明怎样用迹公式研究三角和, 并给出实例.

其余各讲面向的读者较为专门, 文风也相应有所不同. 《\nameref{III}》是上述报告的一个“模”推广, 以 SGA~4 第~XVII~讲 \S5.5 对对称幂的研究为基础. 《\nameref{IV}》在不同情形下定义闭链的上同调类, 并建立相交与杯积的相容性. 《\nameref{V}》汇集了一些已知但缺少文献出处的结果, 以及若干相容性结果. 《\nameref{VII}》是新的. 特别地, 对无支集条件的上同调, 它给出了与紧支集上同调中已知定理类似的有限性定理.

关于各讲的更多说明, 请参阅各讲的引言.

最后, 感谢 J.~L. Verdier 允许我在这里重印他的笔记《\nameref{VIII}》. 我相信这些笔记仍然十分有用, 而它们此前已很难找到.

本卷内部的引文以各讲的简称引用各讲; 这些简称在目录中用方括号标出.

\vspace{20pt}
\begin{flushright}
Bures-sur-Yvette, 1976 年 9 月 20 日\\[5pt]
Pierre Deligne
\end{flushright}
